% =====================================================================
%  PlantillaApuntes.tex
%  ========================
%
%
%  Descripción:
%  ------------
%  Plantilla para los apuntes de clase de cualquier materia. A
%  diferencia de los reportes de actividad, esta es informal y está
%  pensada para llenarse sesión por sesión desde la laptop: se cambian
%  los cuatro datos del bloque DATOS DE LA MATERIA y a partir de ahí
%  solo se agregan bloques \clase.
%
%  Trae entornos listos para definición, teorema, ejemplo y ejercicio,
%  tres cajas de color para lo que se relee antes de un examen (idea
%  clave, cuidado, pendiente) y un estilo de listings para pegar código
%  en C, C++ o Python.
%
%
%  Considerations:
%  ------------
%  - Compila con pdflatex sin shell-escape. Dos pasadas para el índice
%  - babel se carga con es-noshorthands porque en español el " queda
%    activo y se come el texto que le sigue, además de romper listings
%  - Los entornos de teorema comparten contador con las secciones, así
%    que la numeración reinicia en cada clase
%  - Los nombres de los entornos van sin acento (\begin{definicion})
%    aunque el rótulo impreso sí lo lleve
%
%
%  Metadata:
%  ----------
%   * Author : zxxz6 (Bryan Violante Arriaga)
%   * Version: 1.1.0
%
%
%  History:
%  ------------
%  Author      Date            Description
%  zxxz6       18/08/2026      Creation
%
%
% =====================================================================
\documentclass[11pt,letterpaper]{article}

% ==================== IDIOMA Y CODIFICACIÓN ====================
\usepackage[utf8]{inputenc}
\usepackage[T1]{fontenc}
\usepackage[spanish,es-tabla,es-noshorthands]{babel}
\usepackage{lmodern}
\usepackage{microtype}

% ==================== PÁGINA ====================
\usepackage[letterpaper,margin=2.3cm,headheight=14pt]{geometry}
\usepackage{parskip}
\usepackage{fancyhdr}
\usepackage{enumitem}

% ==================== MATEMÁTICAS ====================
\usepackage{amsmath,amssymb,amsthm,mathtools}

% ==================== FIGURAS, TABLAS, CÓDIGO ====================
\usepackage{graphicx}
\usepackage{booktabs}
\usepackage{array}
\usepackage{xcolor}
\usepackage{listings}
\usepackage[most]{tcolorbox}
\usepackage{hyperref}

% =====================================================================
%  DATOS DE LA MATERIA
%  Lo único que hay que cambiar al empezar una libreta nueva.
% =====================================================================
\newcommand{\materia}{Nombre de la materia}
\newcommand{\profesor}{Nombre del profesor}
\newcommand{\periodo}{Otoño 2026}
\newcommand{\alumno}{Bryan Ezequiel Violante Arriaga}

% =====================================================================
%  ESTILO
% =====================================================================

% --- encabezado y pie de cada página ---
\pagestyle{fancy}
\fancyhf{}
\fancyhead[L]{\small\textsc{\materia}}
\fancyhead[R]{\small\periodo}
\fancyfoot[C]{\small\thepage}
\renewcommand{\headrulewidth}{0.4pt}

% --- enlaces ---
\hypersetup{
  colorlinks=true,
  linkcolor=black,
  urlcolor=blue,
  citecolor=black,
  pdftitle={Apuntes de \materia},
  pdfauthor={\alumno}
}

% --- paleta ---
\definecolor{acento}{HTML}{1F4E79}
\definecolor{fondoIdea}{HTML}{EAF2FB}
\definecolor{fondoOjo}{HTML}{FDF0E6}
\definecolor{fondoPend}{HTML}{F2F0F7}
\definecolor{comentario}{HTML}{5A7D5A}
\definecolor{cadena}{HTML}{A03030}
\definecolor{fondoCodigo}{HTML}{F7F7F7}

% --- entornos de teorema; numeran por sección, o sea por clase ---
\theoremstyle{definition}
\newtheorem{definicion}{Definición}[section]
\newtheorem{ejemplo}[definicion]{Ejemplo}
\newtheorem{ejercicio}[definicion]{Ejercicio}

\theoremstyle{plain}
\newtheorem{teorema}[definicion]{Teorema}
\newtheorem{lema}[definicion]{Lema}
\newtheorem{corolario}[definicion]{Corolario}
\newtheorem{proposicion}[definicion]{Proposición}

\theoremstyle{remark}
\newtheorem*{observacion}{Observación}

% --- cajas de color para lo que se relee antes del examen ---
\newtcolorbox{idea}[1][]{
  colback=fondoIdea, colframe=acento, boxrule=0.6pt,
  arc=2pt, left=6pt, right=6pt, top=5pt, bottom=5pt,
  title={Idea clave}, fonttitle=\bfseries\small, #1
}
\newtcolorbox{ojo}[1][]{
  colback=fondoOjo, colframe=orange!70!black, boxrule=0.6pt,
  arc=2pt, left=6pt, right=6pt, top=5pt, bottom=5pt,
  title={Cuidado}, fonttitle=\bfseries\small, #1
}
\newtcolorbox{pendiente}[1][]{
  colback=fondoPend, colframe=violet!60!black, boxrule=0.6pt,
  arc=2pt, left=6pt, right=6pt, top=5pt, bottom=5pt,
  title={Pendiente por resolver}, fonttitle=\bfseries\small, #1
}

% --- código ---
\lstdefinestyle{codigo}{
  backgroundcolor=\color{fondoCodigo},
  basicstyle=\ttfamily\footnotesize,
  keywordstyle=\color{acento}\bfseries,
  commentstyle=\color{comentario}\itshape,
  stringstyle=\color{cadena},
  numbers=left, numberstyle=\tiny\color{gray}, numbersep=8pt,
  frame=single, rulecolor=\color{gray!40},
  breaklines=true, showstringspaces=false,
  tabsize=4, captionpos=b,
  extendedchars=true, inputencoding=utf8,
  literate={á}{{\'a}}1 {é}{{\'e}}1 {í}{{\'i}}1
           {ó}{{\'o}}1 {ú}{{\'u}}1 {ñ}{{\~n}}1
           {Ü}{{\"U}}1 {ü}{{\"u}}1 {¿}{{?`}}1
           {¡}{{!`}}1
}
\lstset{style=codigo}
\renewcommand{\lstlistingname}{Código}
\renewcommand{\lstlistlistingname}{Índice de código}

% --- abreviaturas que se usan a cada rato ---
\newcommand{\N}{\mathbb{N}}
\newcommand{\Z}{\mathbb{Z}}
\newcommand{\Q}{\mathbb{Q}}
\newcommand{\R}{\mathbb{R}}
\newcommand{\C}{\mathbb{C}}
\newcommand{\abs}[1]{\left\lvert #1 \right\rvert}
\newcommand{\norm}[1]{\left\lVert #1 \right\rVert}
\newcommand{\set}[1]{\left\{ #1 \right\}}
\newcommand{\Ord}[1]{\mathcal{O}\!\left(#1\right)}

% --- encabezado de cada sesión: título a la izquierda, fecha a la
%     derecha. El argumento corto es el que sale en el índice ---
\newcommand{\clase}[2]{%
  \section[#1]{#1 \hfill \normalfont\small\textit{#2}}%
}

% =====================================================================
\begin{document}
% =====================================================================

% ==================== PORTADILLA ====================
\begin{center}
  {\LARGE\bfseries Apuntes de \materia\par}
  \vspace{0.4em}
  {\large \profesor \par}
  \vspace{0.2em}
  {\small \periodo \quad$\cdot$\quad \alumno\par}
  \vspace{0.3em}
  \rule{\linewidth}{0.5pt}
\end{center}

\vspace{0.5em}
\tableofcontents
\vspace{1em}
\rule{\linewidth}{0.5pt}

% =====================================================================
%  A PARTIR DE AQUÍ VAN LOS APUNTES.
%  Se copia el bloque de abajo por cada sesión y se borra el contenido
%  de muestra. Todo lo que sigue es un ejemplo de cada herramienta.
% =====================================================================

\clase{Tema de la primera sesión}{12 de agosto de 2026}

Texto normal del apunte. Un párrafo suelto, tal como se escribe en
clase, sin mayor ceremonia. Para resaltar se usa \textbf{negrita} o
\emph{cursiva}, y para nombrar un símbolo en medio de la frase se
escribe matemática en línea: $f(n) = \Ord{n \log n}$.

\begin{definicion}[Nombre de lo que se define]
  Lo que dijo el profesor, redactado con las propias palabras. Los
  entornos se numeran por sesión, así que esta es la 1.1.
\end{definicion}

\begin{teorema}[Nombre del teorema]
  El enunciado. Si $A$ y $B$ son conjuntos finitos, entonces
  \[
    \abs{A \cup B} = \abs{A} + \abs{B} - \abs{A \cap B}.
  \]
\end{teorema}

\begin{proof}
  La demostración, cuando se alcanza a copiar. El rótulo lo pone babel
  en español y el cuadrito del final lo cierra solo.
\end{proof}

\begin{idea}
  Lo que hay que recordar si solo se relee esto antes del examen.
\end{idea}

\begin{ojo}
  El detalle donde siempre se falla. Por ejemplo: el caso base va
  antes del paso inductivo, no después.
\end{ojo}

\begin{pendiente}
  Lo que no se entendió en clase y hay que preguntar o buscar.
\end{pendiente}

\subsection{Un subtema dentro de la sesión}

Ecuación numerada, para poder referirla después:
\begin{equation}\label{eq:ejemplo}
  \sum_{i=1}^{n} i = \frac{n(n+1)}{2}
\end{equation}

En la ecuación~\eqref{eq:ejemplo} la suma es cuadrática en $n$.

Varios pasos alineados por el signo igual:
\begin{align}
  T(n) &= 2\,T\!\left(\frac{n}{2}\right) + n \\
       &= 2\left(2\,T\!\left(\frac{n}{4}\right)
          + \frac{n}{2}\right) + n \nonumber \\
       &= 4\,T\!\left(\frac{n}{4}\right) + 2n
\end{align}

\begin{ejemplo}
  Un ejemplo resuelto en el pizarrón.
  \begin{enumerate}[label=\arabic*., itemsep=0.2em]
    \item Primer paso.
    \item Segundo paso.
  \end{enumerate}
\end{ejemplo}

\begin{ejercicio}
  El que dejó de tarea, para no perderlo entre las hojas.
\end{ejercicio}

\begin{observacion}
  Un comentario al margen, sin número.
\end{observacion}

Código pegado tal cual, con el lenguaje que sea:

\begin{lstlisting}[language=C,caption={Búsqueda binaria}]
int buscar(int arr[], int n, int objetivo) {
    int izq = 0, der = n - 1;
    while (izq <= der) {
        int medio = izq + (der - izq) / 2;   /* evita desbordar */
        if (arr[medio] == objetivo) return medio;
        if (arr[medio] < objetivo) izq = medio + 1;
        else der = medio - 1;
    }
    return -1;
}
\end{lstlisting}

Una tabla, cuando conviene comparar:

\begin{table}[htbp]
  \centering
  \begin{tabular}{lcc}
    \toprule
    \textbf{Algoritmo} & \textbf{Promedio} & \textbf{Peor caso} \\
    \midrule
    Búsqueda lineal   & $\Ord{n}$        & $\Ord{n}$   \\
    Búsqueda binaria  & $\Ord{\log n}$   & $\Ord{\log n}$ \\
    \bottomrule
  \end{tabular}
  \caption{Costo de las dos búsquedas.}
  \label{tab:ejemplo}
\end{table}

% Una figura, cuando se fotografía el pizarrón. Descomentar y poner el
% archivo en la misma carpeta.
%
% \begin{figure}[htbp]
%   \centering
%   \includegraphics[width=0.7\linewidth]{pizarron.png}
%   \caption{Diagrama que dibujó el profesor.}
%   \label{fig:ejemplo}
% \end{figure}

% =====================================================================

\clase{Tema de la segunda sesión}{14 de agosto de 2026}

Se copia este bloque por cada clase nueva y listo.

\end{document}

%######################### END OF PLANTILLAAPUNTES.TEX #########################
%###############################################################################
